\documentclass{article}
%Przykład bazuje na example project overleaf [article]
% Language setting
% Polski
\usepackage[utf8]{inputenc}
\usepackage[MeX]{polski}
\usepackage{indentfirst}
\usepackage{float}
\usepackage{caption}
\usepackage{subcaption}
\usepackage{float}
\usepackage{listings}
\usepackage{amsmath}

% Set page size and margins
% Replace `letterpaper' with`a4paper' for UK/EU standard size
\usepackage[letterpaper,top=2cm,bottom=2cm,left=3cm,right=3cm,marginparwidth=1.75cm]{geometry}

% Useful packages
\usepackage{amsmath}
\usepackage{graphicx}
\usepackage{listings}
\usepackage[]{hyperref}

\usepackage{svg}
\usepackage{dirtree}
\usepackage{algpseudocode}

\usepackage[showframe]{geometry}

\title{Teoria algorytmów i obliczeń - sprawozdanie}
\author{Antoni Adaszkiewicz & Jan Gąska & Jacek Kulma & Mateusz Raczyński}

\begin{document}

\begin{titlepage}
    \newcommand{\HRule}{\rule{\linewidth}{0.5mm}}
	\center
    \includesvg[width=\textwidth,inkscapelatex=false ]{logo-mini.svg}
    \\[0.3cm]
    \HRule\\[0.3cm]
	{\huge\textbf{Teoria algorytmów i obliczeń - sprawozdanie}\\[0.15cm] 
    \HRule\\[2.5cm]
    \begin{center}
        \large
            Antoni Adaszkiewicz \\
            Jan Gąska \\
            Jacek Kulma \\
            Mateusz Raczyński
    \end{center}
	\vfill\vfill\vfill
	{\large \today} \\
	\vfill
\end{titlepage}

\tableofcontents

\newpage
\section{Abstrakt}

\section{Opis problemu}

\section{Rozmiar grafu}

\subsection{Definicja}

\noindent
Przyjęta została następująca definicja \textit{rozmiaru} grafu $G =(V, E)$

\begin{equation}
    |G| = |V| + |E| 
\end{equation}

Definicja ta jest powszechnie stosowana w kontekście złożoności obliczeniowej problemów grafowych (\cite{size}). Użycie rzędu ($|V|$) bądź rozmiaru w kontekście matematycznym ($|E|$) do zdefiniowania \textit{rozmiaru} grafu  oznaczałoby, że graf docelowy może być zwiększany o dowolną liczbę krawędzi/wierzchołków bez zwiększenia rozmiaru (co mija się z intuicją).

\subsection{Algorytm}

\noindent
Poniżej przedstawiony został algorytm do obliczania rozmiaru grafu $G$ reprezentowanego przez macierz sąsiedztwa $M_{G}$:
\\
\begin{algorithmic}
\Function{GetSize}{$M_{G}[\;][\;]$}
    \State {$|V|$ $\gets$ {$GetRowCount(M_{G})$}}
    \State {$Size$ $\gets$ {$|V|$}}
    \For{$i \gets 1$ to $N$}
        \For{$j \gets 1$ to $N$}
            \State {$Size$ $\gets$ {$Size+ M_{G}[i][j]$}}
        \EndFor
    \EndFor
\State \Return {$Sum$}
\end{algorithmic}

\noindent
\\
Złożoność czasowa tego algorytmu to $O(|V|^2)$.

\section{Metryka}
\subsection{Definicja}
\subsection{Algorytm}
\subsection{Algorytm aproksymujący}

\section{Znajdowanie grafu izomorficznego}
\subsection{Definicja}
\subsection{Algorytm}
\subsection{Algorytm aproksymujący}

\section{Testy obliczeniowe}
\subsection{Rozmiar grafu}
\subsection{Metryka}
\subsection{Znajdowanie grafu izomorficznego}

\section{Opis techniczny}
\subsection{Specyfikacja}
\subsection{Budowanie i uruchamianie}

\section{Wnioski}

\renewcommand{\refname}{Bibliografia}
\begin{thebibliography}{9}
\bibitem{size}
\url{https://en.wikipedia.org/wiki/Graph_(discrete_mathematics)}

\end{thebibliography}
\end{document}